\documentclass[12pt,a4paper]{article}

\usepackage[margin=1in]{geometry}
\usepackage{setspace}
\usepackage{graphicx}
\usepackage{booktabs}
\usepackage{amsmath}
\usepackage{amssymb}
\usepackage{siunitx}
\usepackage{float}
\usepackage{hyperref}
\usepackage{caption}
\usepackage{array}

\hypersetup{
    colorlinks=true,
    linkcolor=blue,
    urlcolor=blue,
    citecolor=blue
}

\setstretch{1.2}
\setlength{\parskip}{6pt}
\setlength{\parindent}{0pt}

\title{Case C: Community Microgrid Coursework Report}
\author{Energy Systems Modelling and Engineering}
\date{\today}

\begin{document}

\maketitle

\section{Introduction}

This coursework considers a community microgrid consisting of one shared photovoltaic (PV) system, one shared battery, three household electrical loads, and a single grid connection. The aim is to determine an hourly dispatch schedule that minimises the total electricity cost for the community while respecting battery operating constraints and the physical electricity balance at the community node.

The work follows the modelling framework taught in the course:
\[
\text{Physical system} \rightarrow \text{Mathematical abstraction} \rightarrow \text{Executable code} \rightarrow \text{Verified results}
\]

It also follows the five-layer modelling pipeline:
\[
\text{Data layer} \rightarrow \text{Model layer} \rightarrow \text{Solve layer} \rightarrow \text{Verification layer} \rightarrow \text{Output layer}
\]

The purpose of the model is not only to obtain a least-cost schedule, but also to present a transparent formulation, demonstrate explicit verification, and derive defensible engineering insight from the results.

\section{Dataset and Assumptions}

\subsection{Dataset inspection}

The input dataset is \texttt{caseC\_community\_microgrid\_hourly.csv}. Inspection of the file showed the following columns:

\begin{itemize}
    \item \texttt{timestamp}
    \item \texttt{pv\_kw}
    \item \texttt{load1\_kw}
    \item \texttt{load2\_kw}
    \item \texttt{load3\_kw}
    \item \texttt{import\_tariff\_gbp\_per\_kwh}
    \item \texttt{export\_price\_gbp\_per\_kwh}
\end{itemize}

The total community load is defined by:
\[
L_{\mathrm{tot},t} = L_{1,t} + L_{2,t} + L_{3,t}
\]

The dataset contains 336 rows, corresponding to 14 days of hourly operation. The timestep was confirmed to be \SI{1}{hour} throughout. No missing values were found, no negative values were present in the PV, load, or tariff series, and all input time series had equal length. The units are consistent with the coursework specification: power in \si{kW} and prices in \si{GBP/kWh}.

\subsection{Battery parameters}

The shared battery was modelled using the coursework default parameters:

\begin{table}[H]
\centering
\caption{Battery parameters}
\begin{tabular}{>{\raggedright\arraybackslash}p{8cm}c}
\toprule
Parameter & Value \\
\midrule
Usable capacity, $E_{\max}$ & \SI{10}{kWh} \\
Maximum charge power, $P_{\mathrm{ch,max}}$ & \SI{5}{kW} \\
Maximum discharge power, $P_{\mathrm{dis,max}}$ & \SI{5}{kW} \\
Round-trip efficiency & 90\% \\
Charge efficiency, $\eta_{\mathrm{ch}}$ & 0.948683 \\
Discharge efficiency, $\eta_{\mathrm{dis}}$ & 0.948683 \\
Initial state of charge, $E_0$ & \SI{5}{kWh} \\
Timestep, $\Delta t$ & \SI{1}{h} \\
Grid charging allowed & Yes \\
\bottomrule
\end{tabular}
\end{table}

\subsection{Metering and modelling assumptions}

The metering assumption used in this coursework is:

\emph{A single community meter is located at the grid connection, and all three household loads are aggregated at community level for dispatch optimisation.}

The additional assumptions are:

\begin{itemize}
    \item the system is represented as a single electricity bus
    \item no load shedding is permitted
    \item no explicit battery degradation cost is included
    \item grid import and export are both allowed
    \item the terminal state of charge is constrained to satisfy $E_T = E_0$
    \item PV is split between local use and export through $s_{\mathrm{use},t} + g_{\mathrm{exp},t} = PV_t$
\end{itemize}

The terminal equality condition was chosen so that the optimiser cannot reduce the apparent cost by finishing the horizon with an empty battery.

\section{Mathematical Formulation}

\subsection{Indices}

\[
t = 1,2,\dots,T
\]

where $T=336$.

\subsection{Parameters}

\[
PV_t,\;L_{1,t},\;L_{2,t},\;L_{3,t},\;L_{\mathrm{tot},t},\;\pi_{\mathrm{imp},t},\;\pi_{\mathrm{exp},t},\;\Delta t,\;E_{\max},\;P_{\mathrm{ch,max}},\;P_{\mathrm{dis,max}},\;\eta_{\mathrm{ch}},\;\eta_{\mathrm{dis}},\;E_{\mathrm{init}}
\]

with
\[
L_{\mathrm{tot},t} = L_{1,t} + L_{2,t} + L_{3,t}
\]

\subsection{Decision variables}

\[
s_{\mathrm{use},t},\;p_{\mathrm{ch},t},\;p_{\mathrm{dis},t},\;g_{\mathrm{imp},t},\;g_{\mathrm{exp},t},\;E_t
\]

where:

\begin{itemize}
    \item $s_{\mathrm{use},t}$ is PV used within the community
    \item $p_{\mathrm{ch},t}$ is battery charge power
    \item $p_{\mathrm{dis},t}$ is battery discharge power
    \item $g_{\mathrm{imp},t}$ is grid import power
    \item $g_{\mathrm{exp},t}$ is grid export power
    \item $E_t$ is stored battery energy
\end{itemize}

\subsection{Objective function}

The community-level objective is to minimise the total electricity cost:
\[
\min \sum_{t=1}^{T} \left( \pi_{\mathrm{imp},t} g_{\mathrm{imp},t} - \pi_{\mathrm{exp},t} g_{\mathrm{exp},t} \right)\Delta t
\]

\subsection{Constraints}

Hourly energy balance:
\[
s_{\mathrm{use},t} + p_{\mathrm{dis},t} + g_{\mathrm{imp},t}
=
L_{\mathrm{tot},t} + p_{\mathrm{ch},t} + g_{\mathrm{exp},t}
\]

Battery state update:
\[
E_{t+1} = E_t + \eta_{\mathrm{ch}} p_{\mathrm{ch},t}\Delta t - \frac{1}{\eta_{\mathrm{dis}}}p_{\mathrm{dis},t}\Delta t
\]

State-of-charge bounds:
\[
0 \le E_t \le E_{\max}
\]

Charge and discharge power limits:
\[
0 \le p_{\mathrm{ch},t} \le P_{\mathrm{ch,max}}
\]
\[
0 \le p_{\mathrm{dis},t} \le P_{\mathrm{dis,max}}
\]

PV utilisation and allocation:
\[
0 \le s_{\mathrm{use},t} \le PV_t
\]
\[
s_{\mathrm{use},t} + g_{\mathrm{exp},t} = PV_t
\]

Grid non-negativity:
\[
g_{\mathrm{imp},t} \ge 0,\qquad g_{\mathrm{exp},t} \ge 0
\]

Initial and terminal conditions:
\[
E_0 = E_{\mathrm{init}}
\]
\[
E_T = E_{\mathrm{init}}
\]

\section{Implementation}

The model was implemented in Python using \texttt{pandas}, \texttt{numpy}, \texttt{matplotlib}, and \texttt{cvxpy}. The code follows the same structure as the mathematical formulation and is divided into six stages:

\begin{enumerate}
    \item data loading and inspection
    \item parameter definition
    \item optimisation model construction
    \item solver execution
    \item post-processing and KPI calculation
    \item verification and plotting
\end{enumerate}

The optimisation was solved as a linear programme using \texttt{cvxpy} with the \texttt{CLARABEL} solver. A fallback option to \texttt{SCS} was included in case the preferred solver failed.

\section{Verification}

Verification was treated as an explicit and quantitative stage rather than relying only on solver status.

\subsection{Base-case verification results}

\begin{table}[H]
\centering
\caption{Base-case verification evidence}
\begin{tabular}{lc}
\toprule
Check & Result \\
\midrule
Maximum absolute energy-balance error & \SI{0.000000}{kW} \\
Mean absolute energy-balance error & \SI{0.000000}{kW} \\
Maximum SOC lower-bound violation & \SI{0.000000}{kWh} \\
Maximum SOC upper-bound violation & \SI{0.000000}{kWh} \\
Maximum charge-power violation & \SI{0.000000}{kW} \\
Maximum discharge-power violation & \SI{0.000000}{kW} \\
Maximum PV-limit violation & \SI{0.000000}{kW} \\
Terminal SOC violation & \SI{0.000000}{kWh} \\
Optimisation objective & \pounds 85.020517 \\
Recomputed net cost & \pounds 85.020517 \\
Objective mismatch & \pounds 0.000000 \\
\bottomrule
\end{tabular}
\end{table}

\subsection{Unit consistency}

The unit conversions were checked explicitly:
\[
\si{kW} \times \si{h} = \si{kWh}
\]
\[
\si{GBP/kWh} \times \si{kWh} = \si{GBP}
\]

These checks are important because the optimisation variables are powers, whereas the reported engineering outputs are energies and costs over time.

\section{Results and Discussion}

\subsection{Base-case energy and cost breakdown}

\begin{table}[H]
\centering
\caption{Base-case key performance indicators}
\begin{tabular}{lc}
\toprule
KPI & Value \\
\midrule
Total PV generation & \SI{839.324}{kWh} \\
PV used locally & \SI{784.650}{kWh} \\
PV exported & \SI{54.674}{kWh} \\
Total grid import & \SI{694.033}{kWh} \\
Battery charge energy & \SI{199.590}{kWh} \\
Battery discharge energy & \SI{179.631}{kWh} \\
Total community load & \SI{1404.050}{kWh} \\
Battery losses & \SI{19.959}{kWh} \\
Import cost & \pounds 88.126 \\
Export revenue & \pounds 3.105 \\
Net total cost & \pounds 85.021 \\
Initial SOC & \SI{5.000}{kWh} \\
Final SOC & \SI{5.000}{kWh} \\
SOC minimum & \SI{0.000}{kWh} \\
SOC maximum & \SI{10.000}{kWh} \\
Battery throughput & \SI{379.221}{kWh} \\
\bottomrule
\end{tabular}
\end{table}

The PV self-consumption performance is high. The fraction of PV used locally is:
\[
\frac{784.650}{839.324} \approx 93.5\%
\]

This indicates that the shared battery is effective in absorbing surplus PV and reducing export.

\subsection{Battery value relative to a no-battery reference}

A no-battery reference case was also solved using the same PV, load, and tariff time series. The no-battery case gave:

\begin{itemize}
    \item net total cost = \pounds 119.077
    \item grid import = \SI{784.724}{kWh}
    \item grid export = \SI{109.999}{kWh}
\end{itemize}

Compared with the no-battery reference, the shared battery reduced:

\begin{itemize}
    \item net cost by \pounds 34.056
    \item grid import by \SI{90.691}{kWh}
    \item PV export by \SI{55.325}{kWh}
\end{itemize}

This shows that the battery is primarily valuable for increasing PV self-consumption and avoiding expensive import periods rather than for earning export revenue.

\subsection{Dispatch interpretation}

The battery cycles actively between \SI{0}{kWh} and \SI{10}{kWh}, indicating that the full usable capacity is valuable. Total throughput is \SI{379.221}{kWh} over the 336-hour horizon. The import tariff reaches \pounds 0.3145/kWh, while the export price remains much lower, with a maximum of \pounds 0.0635/kWh. Consequently, the economically preferred strategy is to use stored energy to reduce imports rather than to charge purely for export.

From an engineering perspective, the battery performs three functions:

\begin{itemize}
    \item PV self-consumption improvement
    \item time-shifting of energy across low-value and high-value periods
    \item partial peak shaving during evening load peaks
\end{itemize}

\section{Extension: Feeder Import Limit}

\subsection{Extension formulation}

The selected extension introduces a feeder import limit:
\[
P_{\mathrm{grid,imp,max}} = \SI{4.5}{kW}
\]

The additional constraint is:
\[
0 \le g_{\mathrm{imp},t} \le P_{\mathrm{grid,imp,max}}
\]

This value was chosen as a realistic but binding limit. In the unconstrained base case, the maximum grid import reached \SI{10.686}{kW}, while the maximum net load reached approximately \SI{6.427}{kW}. A \SI{4.5}{kW} limit therefore meaningfully restricts operation but remains feasible when supported by the shared battery.

\subsection{Extension verification and comparison}

\begin{table}[H]
\centering
\caption{Base case and feeder-limit extension comparison}
\begin{tabular}{lcc}
\toprule
Metric & Base case & Extension case \\
\midrule
Net total cost & \pounds 85.021 & \pounds 89.748 \\
Total grid import & \SI{694.033}{kWh} & \SI{690.402}{kWh} \\
Battery throughput & \SI{379.221}{kWh} & \SI{310.230}{kWh} \\
Final SOC & \SI{5.000}{kWh} & \SI{5.000}{kWh} \\
Maximum import violation & -- & \SI{0.000000}{kW} \\
Binding timesteps & -- & 60 \\
\bottomrule
\end{tabular}
\end{table}

The feeder constraint remained feasible with zero unmet demand and zero constraint violations. The import limit was binding in 60 timesteps, concentrated mainly in evening periods around 20:00 to 22:00.

\subsection{Engineering interpretation}

The feeder constraint changes the operational value of the shared battery. In the base case, the battery mainly follows tariff and PV signals. Under the import cap, the battery must also be reserved to protect the community against network import constraints. This increases the total cost by \pounds 4.728 because the battery can no longer be used purely for least-cost arbitrage. The extension therefore shows that shared storage can provide not only market value but also network support value.

\section{Figures}

\begin{figure}[H]
    \centering
    \includegraphics[width=0.92\textwidth]{outputs_caseC/raw_data_sanity_check.png}
    \caption{Raw dataset sanity check showing PV generation and the three household loads.}
\end{figure}

\begin{figure}[H]
    \centering
    \includegraphics[width=0.92\textwidth]{outputs_caseC/base_case_total_load_vs_pv.png}
    \caption{Base case total community load and PV generation. This figure highlights periods of PV surplus and evening net demand.}
\end{figure}

\begin{figure}[H]
    \centering
    \includegraphics[width=0.92\textwidth]{outputs_caseC/base_case_soc.png}
    \caption{Base case battery state of charge. The battery uses its full operating range, indicating that the storage capacity is active and economically relevant.}
\end{figure}

\begin{figure}[H]
    \centering
    \includegraphics[width=0.92\textwidth]{outputs_caseC/base_case_grid_exchange.png}
    \caption{Base case grid import and export. Grid imports are reduced in high-tariff periods through coordinated PV and battery dispatch.}
\end{figure}

\begin{figure}[H]
    \centering
    \includegraphics[width=0.92\textwidth]{outputs_caseC/extension_case_grid_exchange.png}
    \caption{Extension case grid import and export with a \SI{4.5}{kW} feeder import limit. Flat segments at the import ceiling indicate binding network constraints.}
\end{figure}

\section{Conclusion}

This coursework developed, implemented, solved, and verified a linear dispatch model for a community microgrid with shared PV, shared battery storage, three household loads, and grid import/export capability. The model is transparent, reproducible, and directly aligned with the course emphasis on formulation, implementation, verification, and engineering interpretation.

The main findings are:

\begin{itemize}
    \item the battery increased PV self-consumption to approximately 93.5\%
    \item the battery reduced net total cost from \pounds 119.077 in the no-battery reference to \pounds 85.021 in the base case
    \item the battery reduced both grid import and PV export by shifting energy over time
    \item the feeder import limit increased the system cost to \pounds 89.748 and bound in 60 timesteps
    \item under network constraints, the shared battery provides both cost-saving and feeder-support value
\end{itemize}

The main limitations are the simplified single-node representation, the absence of battery degradation cost, and the omission of intra-community allocation or fairness issues. These would be suitable directions for future model extensions.

\section*{AI Usage Statement}

I used generative AI as a support tool to help structure the modelling workflow, improve code readability, and draft report text. I independently checked the dataset interpretation, implemented and ran the optimisation model, reviewed the verification outputs, and edited the final submission to ensure that the modelling assumptions, equations, numerical results, and engineering discussion accurately reflect my own work.

\end{document}
